\documentclass[10pt,a4paper]{amsart}
\usepackage[margin=19mm]{geometry}
\usepackage{amsmath,amssymb,mathtools}
\usepackage[scr=rsfs,cal=euler]{mathalfa}
\usepackage{xfrac}
\usepackage[unicode,hidelinks,bookmarks=false]{hyperref}
\newcommand{\ie}{i.e.\ }
\newcommand{\cf}{cf.\ }
\newcommand{\resp}{respectively\ }
\newcommand{\Subsection}{\subsection}
\providecommand{\nobreakdash}{\nobreak\hbox{-}}
\setlength{\emergencystretch}{2em}
\multlinegap0pt
\usepackage{bm}
\usepackage{bbm}
\usepackage{dsfont}
\usepackage{xspace}
\usepackage{cancel}
\usepackage{mathtools}
\usepackage{amsthm}
\makeatletter
\@ifundefined{thm}{\newtheorem{thm}{Thm}}{}
\@ifundefined{lem}{\newtheorem{lem}[thm]{Lemma}}{}
\@ifundefined{prop}{\newtheorem{prop}[thm]{Proposition}}{}
\makeatother
\def\R{\mathbb R}
\def\eps{\epsilon}
\title[Rado's paracompactness theorem for conformal manifolds]{Rado's paracompactness theorem for conformal manifolds}
\author{Michael Kapovich}
\date{Source: arXiv:2508.02655v1 (CC BY 4.0)}
\begin{document}
\maketitle

\noindent\emph{Source and licence.} Rado's paracompactness theorem for conformal manifolds. Version arXiv:2508.02655v1 (CC BY 4.0), distributed under the Creative Commons Attribution 4.0 International licence, as recorded in the arXiv licence metadata for this exact version and shown on the abstract page https://arxiv.org/abs/2508.02655v1. Full rights notice: licenses/arxiv-2508.02655-CC-BY-4.0.txt.\par\smallskip

\noindent\textit{Editorial scope.} Counted unit: the Prop whose source label is \texttt{None} in the article (source0.tex lines 355--357, source label \texttt{None}), delivered together with the article's own complete original proof of it (source0.tex lines 358--385, one \texttt{proof} environment of 28 lines). No mathematics is written, repaired or re-derived: statement and proof are the article's own text line for line. Required context retained uncounted: lem:L0 (lines 202--206). Every definition, display and label of those lines is retained unchanged. Omitted from this package and not redistributed: 1--201, 207--354, 386--412. Nothing inside a retained excerpt is omitted or altered: every retained body is delivered exactly as the article's lines are written, so no omission receipt of any kind accompanies this package. Citations: the retained excerpts carry no citation command at all, so no bibliography entry of the article is transcribed and no reference is redistributed. Reference adaptation: none was needed and none was made; every reference inside the delivered unit resolves inside it, so no unresolved reference or citation is produced. Printed slips kept literal: no printed word, symbol, hypothesis or number of the counted statement or of its proof was altered. Layout and class: the article's own class is \texttt{article} with packages including geometry, graphicx, amsmath, amsthm, amssymb; the wrapper is the lane's pinned \texttt{amsart} editorial wrapper at 10pt on A4, which restates the theorem-like environments the retained text needs and adds the article's own macro definitions that the retained text needs (\texttt{\textbackslash R}, \texttt{\textbackslash eps}), with extra wrapper packages (bm, bbm, dsfont, xspace, cancel, mathtools). The delivered numbering is the unit's own. Assets: no retained span contains a figure environment, a graphics-inclusion command, a PDF-page inclusion, a table or a TikZ picture; no image file is copied into this package, the package declares no asset dependency and no image file is redistributed with it. Licence: version arXiv:2508.02655v1 of this article is distributed under the Creative Commons Attribution 4.0 International licence, as recorded in the arXiv licence metadata for this exact version and shown on the abstract page https://arxiv.org/abs/2508.02655v1. A full rights notice is delivered at licenses/arxiv-2508.02655-CC-BY-4.0.txt. Characters: every retained excerpt is pure ASCII, so the delivered engine, XeLaTeX with its default Latin Modern text font, reports no missing character.
\begin{lem}\label{lem:L0}
Let $B=B(\mathbf 0, 1)\subset \R^n$ be the open unit ball equipped with some conformal structure (not necessarily the standard one). 
Then for every $\eps>0$ and $r>0$ there exists 
$f\in Lip_c(B)$ which is identically equal to $1$ on the ball $B(\mathbf 0, r)$ and satisfies $I_B(f)<\eps$.  
\end{lem}

\begin{prop}
If $M$ is of Class II, then the metric $\mu_M$ metrizes $M$ as a topological space. 
\end{prop}

\begin{proof} 1. We first prove that the function $\mu_M: M^2\to \R$ is continuous at the diagonal, i.e. if $x_i, y_i$ are sequences in $M$ converging to the same point 
$z\in M$, then $\mu_M(x_i, y_i)\to 0$. Fix an open coordinate  ball $B\subset M$ centered at $z$. Without loss of generality, 
$x_i, y_i\in B$ for all $i$. Then,  
$$
\mu_B(x_i,y_i)\to 0,
$$
see Lemma \ref{lem:L0}. Since $\mu_M(x_i,y_i)\le \mu_B(x_i,y_i)$, we get 
$$
\lim_{i\to\infty} \mu_M(x_i,y_i)= 0= \mu_M(z,z). 
$$ 

2. Let us check continuity of $\mu_M$ at general pairs $(x,y)\in M^2$. Consider sequences $x_i\to x, y_i\to y$ in $M$. 
Then, by the triangle inequality for $\mu_M$ and Part 1 of the proof:
$$
\mu_M(x,y)\le \liminf_{i\to\infty} (\mu_M(x, x_i) + \mu_M(x_i, y_i) + \mu_M(y_i,y))\le   \liminf_{i\to\infty}  \mu_M(x_i, y_i). 
$$ 
Similarly, 
$$
\limsup_{i\to\infty} \mu_M(x_i, y_i)\le \mu_M(x,y). 
$$
It follows that $\mu_M$ is continuous at $(x,y)$. Thus, the manifold topology of $M$ is stronger than the metric topology induced by $\mu_M$. 

3. Let us prove that the metric topology is stronger than the manifold topology. 
Let $B\subset M$ be an open coordinate ball centered at $x\in M$. We will show that there exists $r>0$ such that for every $z\in M- {B}$, $\mu_M(x,z)\ge r$. 
Indeed, since the function $h(y)=\mu_M(x,y)$ is continuous and $S=\partial B$ is compact, the function $h$ has positive minimum $r>0$  on $S$. Take $z\in M- B$ 
and a continuum $C\subset M$ connecting $x$ and $z$. This continuum has to intersect $S$ at some point $y$. Therefore, $r\le \mu(x,y)\le Cap_M(C)$ and, thus,  
$\mu_M(x,z)\ge r$. Proposition follows. 
\end{proof}

\end{document}
